% AI-produced by GPT-6 Astra (OpenAI, ChatGPT Pro mode). No priority claim.
\subsection{Ramsey classes and integral affine equalizers}

The following argument is included with full proof. We do not claim
priority for the implication; the literature search did not identify an
earlier statement in these terms.

\begin{theorem}[Proved: a Ramsey criterion for integral exactness]
\label{OG:Ramsey-exact}
Let $A$ be a small category with finite hom-sets. Suppose that any two
objects have a common successor. Suppose also that $A$ has the following
morphism Ramsey property: for any objects $p,q$ and positive integer $c$
there is an object $r$ such that every coloring
\[
 \chi:A(p,r)\longrightarrow\{1,\ldots,c\}
\]
admits $g:q\to r$ for which $\chi$ is constant on
$g\circ A(p,q)$. Then $A$ has exact colimits in abelian groups.
More precisely, if $a_0,\ldots,a_k:p\to q$ and $k\geq1$, every target
$r$ with the stated Ramsey property for $p,q,c=(3k+1)^k$ admits an
integral affine combination $t\in\mathbb Z[A(q,r)]$ satisfying
$ta_0=\cdots=ta_k$.
\end{theorem}
\begin{proof}
Fix the family and such a target. Put $X=A(p,r)$, $Y=A(q,r)$, and
define the homomorphism between finite-rank free abelian groups
\[
 D:\mathbb Z[Y]\longrightarrow\mathbb Z[X]^k,
 \qquad D[g]=([ga_i]-[ga_0])_{i=1}^k.
\]
Let $\epsilon:\mathbb Z[Y]\to\mathbb Z$ be augmentation, and write
$\epsilon(\ker D)=d\mathbb Z$ with $d\geq0$. We must prove $d=1$.
Suppose otherwise. Choose
\[
 \theta=\begin{cases}
  1/2,&d=0,\\
  \lfloor d/2\rfloor/d,&d\geq2.
 \end{cases}
\]
Then $1/3\leq\theta\leq1/2$ and $d\theta=0$ in the circle group
$\mathbb T=\mathbb R/\mathbb Z$. Consequently
\[
 \phi:\operatorname{im}D\longrightarrow\mathbb T,
 \qquad \phi(Dt)=\theta\epsilon(t)\pmod{\mathbb Z}
\]
is well-defined: two preimages differ by an element of $\ker D$,
whose augmentation is a multiple of $d$.

This homomorphism extends to a homomorphism
$\widetilde\phi:\mathbb Z[X]^k\to\mathbb T$. To see the extension
directly, choose a Smith normal form for the subgroup
$\operatorname{im}D$ of the finite-rank free group. In a suitable basis
$e_1,\ldots,e_m$ it is generated by $s_1e_1,\ldots,s_\ell e_\ell$,
where $s_j>0$. Choose an $s_j$th root in $\mathbb T$ of
$\phi(s_je_j)$ as the value on $e_j$, and assign arbitrary values to
the remaining basis vectors. Such roots exist because multiplication
by every positive integer on $\mathbb T$ is surjective.

Write $\lambda_i(x)\in[0,1)$ for the representative of the value
of $\widetilde\phi$ on the $x$-basis vector in block $i$. For every
$g\in Y$, the construction gives
\begin{equation}\label{OG:circle-obstruction}
 \sum_{i=1}^k\bigl(\lambda_i(ga_i)-\lambda_i(ga_0)\bigr)
       =\theta\pmod{\mathbb Z}.
\end{equation}
Set $L=3k+1$ and color $x\in X$ by the $k$-tuple
\[
 \bigl(\lfloor L\lambda_1(x)\rfloor,\ldots,
       \lfloor L\lambda_k(x)\rfloor\bigr)\in\{0,\ldots,L-1\}^k.
\]
Ramsey supplies $g:q\to r$ on whose $p$-subobjects this color is
constant. In particular, $ga_i$ and $ga_0$ have the same color for
every $i$. Thus
\[
 \left|\sum_{i=1}^k
   \bigl(\lambda_i(ga_i)-\lambda_i(ga_0)\bigr)\right|
 <\frac{k}{3k+1}<\frac13.
\]
Every real representative of $\theta$ modulo $\mathbb Z$ has absolute
value at least $1/3$, contradicting \eqref{OG:circle-obstruction}.
Therefore $d=1$, which means precisely that some $t\in\ker D$ has
augmentation $1$. A family of one arrow requires only $1_q$.
Together with the common-successor hypothesis, simultaneous integral
affine equalization is Isbell--Mitchell's exactness criterion
\cite[Theorem~3.1 and Proposition~3.3]{IsbellMitchell1976}.
\end{proof}

\begin{remark}[Proved: scope of the finite-hom-set hypothesis]\label{OG:finite-hom-scope}
It is used above only to give an elementary Smith normal form proof
of extension to the circle. Divisibility makes $\mathbb T$ an injective
abelian group, so the same argument works without finite hom-sets if
one invokes that standard injectivity theorem. The finite version
already covers the ordered-graph application below.
\end{remark}

\begin{corollary}[Proved: ordered finite graphs are exact]
\label{OG:exact}
Let $\mathcal G$ have one representative of each finite linearly
ordered simple graph as objects and increasing induced graph
embeddings as arrows, composed as functions. Then $\mathcal G$ has
exact colimits in abelian groups. In particular it is not a
counterexample to totality implying exactness.
\end{corollary}
\begin{proof}
The category is small, has finite hom-sets, and has common successors:
take the disjoint union of two graphs and concatenate their orders.
Finite simple graphs with an arbitrary linear order have the Ramsey
property by \cite[Corollary~4.3, printed p.~51]{HubickaNesetril2019}.
The source defines embeddings by preservation and reflection of
relations (printed p.~6), so these are induced embeddings in the graph
case. Because a finite linear order is rigid, an isomorphic induced
copy of $p$ is specified by a unique embedding from $p$; hence its
Ramsey property for copies is exactly the morphism Ramsey property
used in Theorem~\ref{OG:Ramsey-exact}. That theorem applies.
\end{proof}

\subsection{Explicit small ordered-graph certificates}

Here $q$ is the increasing induced path with vertex set
$\{0<1<2\}$ and edge set $\{01,12\}$. A triple $(a,b,c)$ denotes
the embedding of $q$ with images $a<b<c$. These examples were first
found by exact finite computation, then checked by the displayed
integer identities. They are independent of the Ramsey theorem.

\begin{proposition}[Proved: three vertex maps]
\label{OG:vertex-certificate}
For the three embeddings $v_0,v_1,v_2$ of the one-vertex graph into
$q$, an integral affine equalizer exists in the ordered graph $r_6$
on $\{0<\cdots<5\}$ with edge set
\[
 \{01,12,13,15,23,24,34,35\}.
\]
It is
\[
 t=(012)-(015)+(124)-(134)+(235).
\]
\end{proposition}
\begin{proof}
Each displayed triple has its consecutive pairs in the edge set and
its outer pair absent, so it is an increasing induced embedding.
The coefficient sum is $1-1+1-1+1=1$. The three vertex marginals are
respectively
\[
 [0]-[0]+[1]-[1]+[2]=[2],\quad
 [1]-[1]+[2]-[3]+[3]=[2],\quad
 [2]-[5]+[4]-[4]+[5]=[2].
\]
Thus $tv_0=tv_1=tv_2$ over $\mathbb Z$, and after change of
coefficients over every unital coefficient ring.
\end{proof}

\begin{proposition}[Proved: two edge maps]
\label{OG:edge-certificate}
For the embeddings $e_0,e_1$ of the ordered edge into the two edges
of $q$, an integral affine equalizer exists in the ordered graph $r_7$
on $\{0<\cdots<6\}$ with edge set
\[
 \{01,12,23,35,56,14,45\}.
\]
It is
\[
 t=(012)+(123)+(235)+(356)-(014)-(145)-(456).
\]
\end{proposition}
\begin{proof}
Each displayed triple is an increasing induced path by inspection of
the specified edge set. The augmentation is $4-3=1$. In the free
abelian group on embeddings of the edge, one has
\[
\begin{aligned}
 (e_1-e_0)_*\bigl((012)+(123)+(235)+(356)\bigr)
   &=[56]-[01],\\
 (e_1-e_0)_*\bigl((014)+(145)+(456)\bigr)
   &=[56]-[01].
\end{aligned}
\]
Indeed the two telescoping sums are respectively
\[
\begin{gathered}
 ([12]-[01])+([23]-[12])+([35]-[23])+([56]-[35]),\\
 ([14]-[01])+([45]-[14])+([56]-[45]).
\end{gathered}
\]
Subtracting gives $te_0=te_1$. Precomposition with each of the two
vertex maps into the edge also gives $tv_0=tv_1=tv_2$.
\end{proof}

\paragraph{Finite computations and their limits.}
The supplied script \texttt{finite\_ordered\_graphs.py} enumerates
every graph on the fixed ordered vertex set $\{0,\ldots,n-1\}$ for
$1\leq n\leq6$. It uses exact rational and finite-field arithmetic
to compare the rank of $D$ with the rank after adjoining augmentation
as an extra row. The counts of targets admitting affine equalizers
are zero for both families over all three fields when $n\leq5$.
For $n=6$, among all $2^{15}=32768$ ordered graphs, the counts are
\[
\begin{array}{c|rrr}
 &\mathbb Q&\mathbb F_2&\mathbb F_3\\\hline
 \text{three vertex maps}&42&34&42\\
 \text{two edge maps}&8&0&8
\end{array}
\]
The first vertex certificate has ranks $5,6$ over all three fields;
the displayed seven-vertex edge certificate has ranks $6,7$ over all
three. A separate script, \texttt{finite\_ordered\_graphs\_verify.py},
rebuilds the graph maps directly, checks all reported coefficients
and marginals, and independently computes these ranks using
\texttt{DomainMatrix} over the three exact coefficient domains.
The exhaustive computation shows that the smallest target size for
the vertex family over $\mathbb Z$ is $6$, and for the edge family
over $\mathbb Z$ is $7$: an integral certificate reduces modulo $2$,
so the enumerated $\mathbb F_2$ obstructions give the latter lower
bound. These minimal-size assertions are certified finite
computations, not conceptual lower-bound proofs.

A separately recorded sample of $20000$ seven-vertex graphs, using
seed $20261113$, is expressly not an exhaustive seven-vertex search.
None of these finite searches by itself proves universal exactness
or totality. The universal conclusion in Corollary~\ref{OG:exact}
instead follows from the circle-group argument and the opened
structural Ramsey theorem.
